% ============================================================
%  00_tom_tat.tex
% ============================================================

\section*{Tóm tắt nội dung}
\addcontentsline{toc}{section}{Tóm tắt nội dung}

Bài thực hành số 03 thuộc học phần CSE703093 -- An toàn phần mềm tập trung vào
ba chủ đề cốt lõi: phát hiện lỗi bộ nhớ bằng phân tích động, sửa lỗi theo các
nguyên tắc lập trình C an toàn, và định lượng mức độ rủi ro của từng lỗ hổng
theo mô hình đơn giản hóa theo tinh thần CVSS.

Module thực hành được xây dựng trong bối cảnh hệ thống SecureSys -- Inventory
(quản lý kho sách thư viện). File \textit{c/vulnerable\_inventory.c} cố ý chứa
ba lỗi bộ nhớ kinh điển: Buffer Overflow (CWE-121) thông qua
\textit{strcpy()} không giới hạn độ dài, Use-After-Free (CWE-416) khi
tiếp tục đọc con trỏ sau khi đã \textit{free()}, và Memory Leak
(CWE-401) do thiếu \textit{free()} trong vòng lặp cấp phát. File
\textit{c/secure\_inventory.c} là bản đã sửa, trong đó \textit{strcpy} được
thay bằng \textit{snprintf}, con trỏ được đặt về \textit{NULL} ngay sau
\textit{free}, và mọi vùng nhớ cấp phát đều được giải phóng.

Công cụ AddressSanitizer (ASan) được sử dụng để phát hiện lỗi tại thời
điểm chạy. Kết quả thực nghiệm: ASan phát hiện đủ 3/3 lỗi trên bản vulnerable
(\textit{heap-buffer-overflow}, \textit{heap-use-after-free}, memory leaked by
LeakSanitizer), trong khi bản secure chạy sạch 0/3 lỗi.

Script Python \textit{python/risk\_assessor.py} áp dụng mô hình CVSS-lite để
tính điểm rủi ro cho năm lỗ hổng. Kết quả: VULN-001 (CWE-121) và VULN-005
(CWE-134) cùng đạt 9.8/CRITICAL; VULN-002 (CWE-416) đạt 6.9/MEDIUM; VULN-003
(CWE-401) đạt 5.4/MEDIUM; VULN-004 (CWE-476) đạt 2.5/LOW.

Bộ kiểm thử tự động \textit{tests/test\_memory\_safety.py} gồm 8 bài kiểm tra
và đạt kết quả 8/8 PASS: 3 bài xác nhận ASan phát hiện lỗi, 3 bài xác nhận
bản secure sạch, 1 bài kiểm tra tính nhất quán với baseline, 1 bài xác nhận
VULN-001 là lỗ hổng CRITICAL cao nhất.

Điểm sư phạm quan trọng: tràn bộ đệm \emph{bên trong} một struct (ghi từ trường
\textit{title} sang \textit{author}) có thể không bị ASan phát hiện vì ASan chỉ
đặt redzone ở biên vùng cấp phát, không đặt giữa các field của cùng một
allocation. Đây là lý do vì sao ``chạy không báo lỗi'' không đồng nghĩa với
``chương trình an toàn''.

% ────────────────────────────────────────────────────────────
\section*{Mục tiêu học tập}
\addcontentsline{toc}{section}{Mục tiêu học tập}

\begin{itemize}
    \item Nhận diện, tái hiện và giải thích cơ chế của ba lớp lỗi bộ nhớ phổ
          biến trong C: Buffer Overflow (CWE-121), Use-After-Free (CWE-416) và
          Memory Leak (CWE-401).
    \item Sử dụng AddressSanitizer (\textit{-fsanitize=address}) kết hợp
          LeakSanitizer (\textit{ASAN\_OPTIONS=detect\_leaks=1}) để phát hiện lỗi
          tại runtime và đọc hiểu thông báo lỗi, stack trace.
    \item Sửa lỗi theo nguyên tắc lập trình C an toàn: dùng \textit{snprintf}
          thay \textit{strcpy}, đặt con trỏ về \textit{NULL} sau \textit{free},
          đảm bảo mọi \textit{malloc} có \textit{free} tương ứng, kiểm tra
          giá trị trả về của \textit{malloc}.
    \item Áp dụng mô hình đánh giá rủi ro định lượng theo tinh thần CVSS để
          xếp hạng mức độ nghiêm trọng và ưu tiên thứ tự khắc phục.
    \item Hiểu rõ giới hạn của ASan: tràn bộ đệm trong cùng một struct
          allocation không nhất thiết bị phát hiện, vì redzone chỉ bao quanh
          biên của vùng cấp phát, không phân tách các trường bên trong.
    \item Thực hành viết và chạy bộ kiểm thử tự động tích hợp cả phát hiện lỗi
          runtime và kiểm tra tính đúng đắn của mô hình đánh giá rủi ro.
\end{itemize}

% ────────────────────────────────────────────────────────────
\section*{Yêu cầu hoàn thiện}
\addcontentsline{toc}{section}{Yêu cầu hoàn thiện}

\begin{enumerate}
    \item Bảng lỗi và stack trace: Liệt kê đủ 3 lỗi phát hiện bằng ASan
          (loại lỗi, hàm gây ra, stack trace rút gọn, hậu quả).
    \item Giải thích tràn trong struct: Mô tả và chứng minh bằng thực
          nghiệm tại sao tràn từ \textit{title} sang \textit{author} (cùng một
          struct trên stack) không bị ASan báo lỗi, trong khi tràn heap buffer
          riêng bị bắt ngay lập tức.
    \item Bản vá code: Trình bày giải thích từng sửa đổi trong
          \textit{secure\_inventory.c} tương ứng với từng lỗi.
    \item Báo cáo đánh giá rủi ro đầy đủ: Bảng 5 lỗ hổng (bao gồm
          VULN-005 tự thêm) với điểm số, mức độ và thứ tự ưu tiên khắc phục.
    \item Log kiểm thử: Toàn bộ 8 test của
          \textit{test\_memory\_safety.py} đều PASS.
\end{enumerate}

% ────────────────────────────────────────────────────────────
\section*{Yêu cầu kết quả}
\addcontentsline{toc}{section}{Yêu cầu kết quả}

\begin{itemize}
    \item ASan phát hiện đủ 3/3 lỗi trên bản \textit{vulnerable}; xác nhận 0/3
          lỗi trên bản \textit{secure}.
    \item \textit{tests/test\_memory\_safety.py} toàn bộ 8 test PASS.
    \item Kết quả đánh giá rủi ro khớp với
          \textit{dataset/expected\_risk\_baseline.json} với dung sai tối đa
          0.05 điểm.
\end{itemize}
